\section{Registros de decisiones de arquitectura (ADR)}
\subsection{ADR-001: estilo de arquitectura}
La estructura del registro sigue la guía de AWS~\cite{aws-adr}. Su estado es \textbf{Propuesto}. Se comparan tres alternativas: conservar la provisión de datos académicos dentro del backend modular, integrar el núcleo con un SGA simulado mediante una API, o dividir también responsabilidades internas en microservicios. La mensajería se considera un complemento para necesidades que admitan procesamiento diferido, no un requisito de la consulta de E1. Una entrega al menos una vez puede repetir el procesamiento, como ilustra el caso de Apache Storm~\cite{bonte2025-stream}. La segunda alternativa se representa en la Figura~\ref{fig:c4-contenedores}: una API modular del SGRA consulta al SGA simulado sin acceder directamente a sus datos internos. Este cotejo documental no aprueba el ADR ni establece replicación. Archivo fuente: \path{docs/adr/ADR-001-estilo-arquitectura.md}.
\subsection{ADR-002: mecanismo de comunicación}
La estructura del registro sigue la guía de AWS~\cite{aws-adr}. Su estado es \textbf{Propuesto}. Se recomienda REST/HTTP JSON síncrono para que un servicio de refuerzos consulte una asignatura sintética a un servicio académico, ambos como procesos independientes. Se comparan gRPC y mensajería por contrato, herramientas y dependencia temporal, sin atribuir resultados de rendimiento no medidos. La propuesta define rutas, respuestas \texttt{200/400/404/503}, una sola tentativa y un presupuesto inicial de dos segundos para la llamada saliente. Esta configuración no es una medición ni un RNF validado. La demo local empleará datos sintéticos en memoria y escucha exclusiva en loopback, sin autenticación por sesión. El C4 representa la interacción académica propuesta; faltan las mediciones de PE-U1, la comprobación del esqueleto ejecutado y la revisión del equipo. El contrato y los criterios de aceptación están en \path{docs/adr/ADR-002-mecanismo-comunicacion.md}. La propuesta no cambia la decisión de estilo arquitectónico de ADR-001.
\subsection{ADR-003: alcance, reutilización, tecnologías y datos críticos}
Su estado es \textbf{Aprobado}. El alcance comprende a estudiantes, docentes, coordinación académica, administración y el gestor de áreas de trabajo. Coordinación realiza seguimiento; en el diseño propuesto, la información académica se consulta mediante la API del SGA simulado. Groq se propone como apoyo exclusivo a la validación del formulario de solicitudes de refuerzo. Se reutiliza la organización por capas y dominios del backend y la estructura de la SPA por roles. La base inicial del backend contemplaba Java 21, Spring Boot 4.0.3 y Maven; el receptor académico de E1 se generó con Spring Boot 4.0.10, ofrecido por Spring Initializr, manteniendo Java 21 y Maven. La Figura~\ref{fig:c4-contenedores} distingue ambos alcances. La Tabla~\ref{tab:cap-datos} recoge los grupos de datos críticos de este ADR y los datos del SGA que intervienen en la validación académica. Archivo fuente: \path{docs/adr/ADR-003-alcance-reutilizacion-tecnologias.md}.
